\begin{pastexam}[label={px:ntu-111-bime-ds-02}]{臺大生機所 111 學年度（C）第 2 題}
（問答）

(15\%) Given a graph, write a pseudocode function that colors the graph's vertices such that no
two adjacent vertices share the same color by minimizing the total number of colors used.
\end{pastexam}

\begin{pastsoln}
用的顏色越少越好，就是求圖的色數（資結書第~20.1.4~節）。 做法：依序試 $m = 1, 2, 3, \ldots$，每個
$m$ 用回溯檢查「能不能用 $m$ 種顏色塗」，第一個成功的 $m$ 就是最少的顏色數。 回溯就是資結書
虛擬碼~20.4 的 $m$ 著色，改成找到第一種塗法就停。 頂點編號 $1$ 到 $n$（$n \ge 1$），$x_i$ 是頂點 $i$
的顏色，$\mathit{found}$ 是全域的旗標。

\begin{pseudo}{用最少顏色著色}
\Procedure{Try-Color}{$k, m$} \Comment{頂點 $1..k-1$ 已經塗好}
    \For{$c \gets 1$ \To $m$}
        \If{$k$ 的每個鄰居 $j < k$ 都有 $x_j \ne c$}
            \State $x_k \gets c$
            \If{$k = n$}
                \State $\mathit{found} \gets \TRUE$
            \Else
                \State \Call{Try-Color}{$k + 1, m$}
            \EndIf
            \If{$\mathit{found}$}
                \State \Return \Comment{找到一種塗法就不再試別的顏色}
            \EndIf
        \EndIf
    \EndFor
\EndProcedure
\Statex
\Function{Min-Coloring}{$G$}
    \For{$m \gets 1$ \To $n$}
        \State $\mathit{found} \gets \FALSE$
        \State \Call{Try-Color}{$1, m$}
        \If{$\mathit{found}$}
            \State \Return $m$ 和塗法 $(x_1, \ldots, x_n)$
        \EndIf
    \EndFor
\EndFunction
\end{pseudo}

\textbf{正確性。} \proc{Try-Color}$(1, m)$ 走的是 $m$ 著色的狀態空間樹，只剪掉「新塗的頂點和已經塗
過的鄰居同色」的節點，這種節點不管後面怎麼塗都有一條邊兩端同色，子樹裡沒有合法的塗法。 其他節點
都會被走到（除非已經找到），所以 $\mathit{found}$ 變成 \TRUE{} 若且唯若圖能用 $m$ 種顏色塗，而且那時
$x_1, \ldots, x_n$ 就是一種合法的塗法。 $m$ 由小到大試，第一個成功的 $m$ 之前每一個都不能塗，所以它
就是最少的顏色數。 $m = n$ 時每個頂點一種顏色一定可以，所以迴圈一定會回傳。

\textbf{時間。} 對固定的 $m$，狀態空間樹最多 $\sum_{i=0}^{n} m^i = O(m^n)$ 個節點，用鄰接串列檢查
\textbf{if} 的條件，每個節點 $O(n)$，所以一次 \proc{Try-Color}$(1, m)$ 是 $O(n m^n)$。 設最少要 $\chi$
種顏色，總時間是 $\sum_{m=1}^{\chi} O(n m^n) = O(n \chi^{n+1})$，最壞是指數時間。 這個問題的決策版本在
$m \ge 3$ 時是 NP 完備的（資結書第~28~章），目前沒有已知的多項式時間演算法；實務上可以固定頂點 $1$
塗顏色 $1$ 去掉對稱，或先塗鄰居多的頂點，讓搜尋少走一些節點（資結書第~20.1.4~節）。

\textbf{不要寫成單純的貪心。} 「每個頂點依序塗上鄰居沒用過的最小顏色」只要 $O(n + m_E)$（$m_E$ 是
邊數），塗出來一定合法，但顏色不一定最少。 例如四個頂點、邊 $\{1, 3\}$、$\{3, 4\}$、$\{4, 2\}$（一條
路徑 $1 - 3 - 4 - 2$）：依編號塗，頂點 1、2 都塗 $1$，頂點 3 塗 $2$，頂點 4 的鄰居用了 $1$、$2$，只好塗
$3$；其實兩種顏色就夠（$1$、$4$ 一色，$2$、$3$ 另一色）。 題目要求最少，所以要用上面的搜尋；用
\proc{Min-Coloring} 算這個圖，$m = 2$ 就成功。
\end{pastsoln}
